% =====================================================================
\section{Interventions: Steering Generation and Stress-Testing Coverage}
\label{sec:interv}

A correlational principle should be stressed interventionally. We probe three
questions on a procedurally generated grid (Kubric-derived motion interventions
\cite{greff2021kubric} plus our SDF-fractal generator), constructed after the
analysis of Secs.~\ref{sec:law}--\ref{sec:absolute} was frozen: \emph{(i)}
construct validity, \emph{(ii)} whether a tuned source beats a generic baseline,
and \emph{(iii)} whether motion (not appearance) is the target-specific carrier
and the coverage response has the predicted shape.

\begin{table}[t]
\centering
\caption{\textbf{Recovered generator parameters $\theta$} from BFV-only
search. The camera/object decomposition is read directly off $\theta$ and
matches each target's known physics; full parameter set in supplement.}
\label{tab:recovery}
\small
\begin{tabular}{l cc l}
\toprule
Target & camera dolly $\Delta$ & obj.\ vert.\ frac. & recovered regime \\
\midrule
KITTI-2015 & \textbf{2.18} & \textbf{0.00} & camera-centric \\
FlyingThings & \textbf{0.32} & \textbf{0.92} & object-centric \\
\bottomrule
\end{tabular}

\end{table}

\noindent\textbf{(i) The recovered motion is physically plausible
(construct validity).} A descriptor as crude as a 4-D bag of flow vectors
invites a natural worry: does it capture anything real? We test this directly.
We run a render-free search over generator motion parameters, scored only by how
well the generated motion covers the target's, and ask what motion matching the
descriptor actually produces (\emph{render-free}: each candidate costs only a
geometry-only flow pass---the scene's motion field, with no texture or lighting
render---and no training; only the winning source is later materialized in full,
with appearance, to train on and verify in (ii)).
On a
controllable Kubric-style generator whose parameters $\theta$ explicitly
separate camera motion from object motion, optimizing BFV match alone recovers
the expected motion regime for each target (Table~\ref{tab:recovery}): a strong
forward dolly with grounded objects for
KITTI (camera-centric), a near-static camera with vertically flying objects for
FlyingThings (object-centric). The recovered camera axis cleanly separates the
moving-camera target from the static-camera target, with no supervision beyond
the motion histogram. The recovered KITTI source's flow fingerprint bears this
out, a coherent radial-expansion field that captures the target's
forward-ego-motion structure, unlike the isotropic generic anchor
(Fig.~\ref{fig:intervsplats}). An independent generator, matched on BFV alone,
lands on the right physical motion: construct validity earned, not asserted---and
not circular, since the objective fixes only the aggregate motion histogram, yet
the regime it selects is the physically correct one, with its real test in (ii). One
caveat: BFV is by construction blind to the camera-vs-object \emph{decomposition}
when the two produce identical aggregate flow, since an ego-motion dolly and
flying objects can share an aggregate fingerprint. We therefore read the
camera/object split off $\theta$, not off the splat, and treat the ambiguity in
the supplement.

\begin{figure*}[t]
\centering
\includegraphics[width=\linewidth]{figures/results/F_pipeline_closedloop.png}
\caption{\textbf{Closed-loop, render-free source design.} A generic source
(MOVi-F) is tuned by a TPE search whose only objective is BFV coverage ($-\dbt$)
against the target, scored from geometry-only flow passes with no texture render
and no training. The recovered ``KITTI-recovered'' source has a coherent
forward-dolly fingerprint ($|f|{\approx}13$) vs.\ the near-isotropic generic
anchor ($|f|{\approx}4$); training on it beats MOVi-F across every architecture
(Table~\ref{tab:interv}). Panels are directional flow splats; color $=$ flow
direction (wheel in supp.).}
\label{fig:intervsplats}
\end{figure*}

\begin{table*}[t]
\centering
\caption{\textbf{A target-tuned source beats the generic baseline.} Peak \pck{}
of the render-free \emph{tuned} source (the KITTI-recovered generator,
GSO assets $+$ HDRI) vs.\ the generic MOVi-F baseline \cite{greff2021kubric},
by architecture. $\Delta =$ tuned $-$ MOVi-F; bold marks a win. The tuned source
wins all eight KITTI cells. On the semantic \textsc{tss} benchmark, the
target-matched source is instead \emph{near-static}, and it too beats MOVi-F on
all three architectures (RAFT, an optical-flow network, is not run on the
semantic benchmarks). The full grid including the from-scratch backbone
configurations---an out-of-distribution probe, where the tuned source still wins
by prior-matching rather than learned features---is reported in full in the
supplement.}
% TODO(DAVIS): describe the TAP-Vid-DAVIS columns here once the intervention run
%              lands; cells are \tbd placeholders for now.
\label{tab:interv}
\setlength{\tabcolsep}{5pt}
\resizebox{\textwidth}{!}{\begin{tabular}{ll rrrrrr}
\toprule
 &  & \multicolumn{3}{c}{KITTI-2012} & \multicolumn{3}{c}{KITTI-2015} \\
\cmidrule(lr){3-5}\cmidrule(lr){6-8}
Model & Source & @5\% & @3\% & @1\% & @5\% & @3\% & @1\% \\
\midrule
\multicolumn{8}{l}{\emph{Kubric generator}} \\
\addlinespace[1pt]
\multirow{2}{*}{CATs++} & MOVi-F & 95.6 & 89.1 & \textbf{45.5} & 94.9 & \textbf{88.6} & \textbf{44.4} \\
 & tuned & \textbf{98.2} & \textbf{89.1} & 37.8 & \textbf{96.5} & 86.0 & 37.0 \\
\addlinespace[2pt]
\multirow{2}{*}{GLU-Net} & MOVi-F & 79.8 & 52.8 & 15.5 & 66.9 & 49.0 & 22.6 \\
 & tuned & \textbf{96.2} & \textbf{85.8} & \textbf{46.8} & \textbf{77.6} & \textbf{60.5} & \textbf{30.7} \\
\addlinespace[2pt]
\multirow{2}{*}{FlowFormer} & MOVi-F & 75.9 & 48.5 & 13.8 & 63.1 & 46.2 & 20.8 \\
 & tuned & \textbf{94.6} & \textbf{80.4} & \textbf{29.5} & \textbf{73.2} & \textbf{56.0} & \textbf{25.8} \\
\addlinespace[2pt]
\multirow{2}{*}{RAFT} & MOVi-F & 84.5 & 62.2 & 25.9 & 74.4 & 57.3 & 30.2 \\
 & tuned & \textbf{96.6} & \textbf{87.1} & \textbf{49.4} & \textbf{78.8} & \textbf{61.8} & \textbf{32.4} \\
\midrule
\multicolumn{8}{l}{\emph{SDF-Fractals3D generator}} \\
\addlinespace[1pt]
\multirow{2}{*}{CATs++} & default & 97.2 & 93.8 & 40.2 & 96.1 & 92.1 & 44.5 \\
 & tuned & \textbf{98.7} & \textbf{97.4} & \textbf{89.1} & \textbf{98.0} & \textbf{95.9} & \textbf{76.3} \\
\addlinespace[2pt]
\multirow{2}{*}{GLU-Net} & default & 89.0 & 74.2 & 22.3 & 79.7 & 69.2 & 26.5 \\
 & tuned & \textbf{98.0} & \textbf{96.2} & \textbf{89.1} & \textbf{94.8} & \textbf{91.8} & \textbf{81.0} \\
\addlinespace[2pt]
\multirow{2}{*}{FlowFormer} & default & 89.7 & 54.4 & 6.2 & 83.7 & 56.7 & 9.7 \\
 & tuned & \textbf{96.8} & \textbf{91.9} & \textbf{62.7} & \textbf{91.2} & \textbf{84.2} & \textbf{56.0} \\
\bottomrule
\end{tabular}
}
\end{table*}

\noindent\textbf{(ii) A tuned source beats a generic baseline.}
The recovered motion also trains well. Materializing the KITTI-recovered source
(GSO assets, HDRI lighting) and training on it beats MOVi-F~\cite{greff2021kubric}
in \emph{all eight} KITTI settings, with gains of $+1.8$ to $+22.5$ \pck{}
across the four architectures (Table~\ref{tab:interv}). A target-specific source
built without a single real-target label beats a generic synthetic dataset. (The
from-scratch backbone configurations, where the tuned source still wins by
prior-matching rather than learned features---a source whose motion marginal sits
closer to the target transfers better even without a trained encoder---are an
out-of-distribution probe deferred to the supplement.)

\noindent\textbf{(iii) Motion is the carrier, and the coverage response is
inverted-U.} Scaling camera-dolly magnitude across five rungs ($0.25\times$ to
$2\times$ KITTI's) and crossing with four texture variants
(Fig.~\ref{fig:kubric_appearance}), we train GLU-Net from
scratch on each of the $20$ sources and read peak \pck{} on KITTI
(Fig.~\ref{fig:ladder}). Transfer is \emph{inverted-U} in magnitude, peaking at
the matched rung for both KITTI splits and every texture variant; missing the
match costs $5$--$11$ \pck. Off-target mass ($\dtb$) stays flat across the ladder
while coverage ($\dbt$) tracks the PCK curve directly---coverage is the operative
lever. Because magnitude moves both directions together, this maps the coverage
\emph{response} rather than isolating it; the flat $\dtb$ nonetheless indicates
the off-target direction is not what drives the curve, and a clean fixed-off-target
isolation remains future work. Appearance leaves the curve's shape unchanged, with no
consistent edge for higher visual fidelity (Fig.~\ref{fig:ladder}), confirming
motion as the target-specific carrier. (The one exception is scratch robustness
to a gross over-shoot: matte/$2\times$ falls to $82.9$ on KITTI-2012
vs.\ photoreal's $91.6$, but a pretrained backbone restores the tolerance; supp.)

\begin{figure}[t]
\centering
\begin{subfigure}[t]{0.37\linewidth}
  \centering
  \includegraphics[height=3.0cm]{figures/results/ladder_panel_a.png}
  \caption{Directed motion distance ($\dbt$, $\dtb$) vs.\ source magnitude.}
  \label{fig:ladder:coverage}
\end{subfigure}%
\hfill%
\begin{subfigure}[t]{0.61\linewidth}
  \centering
  \includegraphics[height=3.0cm]{figures/results/ladder_panel_bc.png}
  \caption{Transfer to KITTI-2012 and KITTI-2015 (peak \pck{}).}
  \label{fig:ladder:heatmaps}
\end{subfigure}

\vspace{3pt}
\begin{subfigure}[b]{\linewidth}
  \centering
  \includegraphics[width=\linewidth]{figures/results/ladder_panel_d.png}
  \caption{Coverage gap ($\dbt$) resolved across the image frame, per source magnitude.}
  \label{fig:ladder:spatial}
\end{subfigure}
\caption{\textbf{Controlled motion $\times$ appearance grid $\to$ KITTI transfer.}
Columns scale source motion magnitude purely to sweep coverage; rows vary appearance.
\emph{(a)} Coverage ($\dbt$, blue; lower${=}$better) dips at the match and grows on
either side, while off-target mass ($\dtb$, red) stays flat.
\emph{(b)} \pck{} varies mostly with coverage (columns) and little with appearance
(rows): it peaks at the match and falls on either side, following coverage (cf.\ (a)),
while higher visual fidelity gives no consistent edge (HDRI+GSO does not win cleanly;
matte+KuBasic does well) despite the rows' DINO appearance-coverage gap (\textbf{bold})
varying about $2\times$. Appearance is largely decoupled from transfer.
\emph{(c)} The same $\dbt$ across the KITTI frame, each region relative to the best
coverage it reaches across the sweep (green${=}$that best, red${=}$worse; splits
pooled): the match is best frame-wide and degrades on either side, most in the
high-motion lower field. Per-model and variant breakdown in the supplement.}
\label{fig:ladder}
\end{figure}

\begin{figure}[t]
\centering
\includegraphics[width=\linewidth]{figures/results/gallery_kubric_appearance.png}
\caption{\textbf{Kubric appearance variants} crossed in the motion grid
(Fig.~\ref{fig:ladder}), at bit-identical motion. Two
knobs: background/lighting---an HDRI environment with textured materials vs.\ a
matte background with flat shading---and object assets---KuBasic procedural
primitives vs.\ GSO (Google Scanned Objects). The DINO appearance distance to
KITTI is driven by the background/lighting knob (HDRI ${\approx}0.6$ vs.\ matte
${\approx}1.1$), not the asset knob (KuBasic ${\approx}$ GSO).}
\label{fig:kubric_appearance}
\end{figure}


